\begin{proof}[Proof of \cref{obj:lem:armwise-modulus}]
\emph{1. Mass and denominator perturbations.}
Use the quantities in \cref{obj:def:armwise-extension}, and define
\[
\delta:=\sum_{a=0}^1\Delta_a.
\]
Each \(\Delta_a\) is nonnegative, so \(\delta\geq0\). The triangle inequality gives, for each \(a\in\{0,1\}\),
\[
\begin{aligned}
|s_a(u)-s_a(v)|
&=\bigl|(u_{a0}-v_{a0})+(u_{a1}-v_{a1})\bigr|
\leq\Delta_a,\\
|s(u)-s(v)|
&\leq |s_0(u)-s_0(v)|+|s_1(u)-s_1(v)|
\leq\delta.
\end{aligned}
\]
Because the coordinates of \(v\) are nonnegative, \(s(v)\geq0\); equality would force every coordinate to vanish. Thus \(v\neq0\) implies
\[
s(v)>0,\qquad
D_a(v)\geq\epsilon s(v)>0,\qquad
\frac{s(v)}{D_a(v)}\geq0.
\]
Comparing the two entries of each maximum yields
\[
\begin{aligned}
|D_a(u)-D_a(v)|
&\leq
\max\bigl\{|s_a(u)-s_a(v)|,\,
          |\epsilon s(u)-\epsilon s(v)|\bigr\}\\
&\leq\max\{\Delta_a,\epsilon\delta\}
\leq\Delta_a+\epsilon\delta.
\end{aligned}
\]
Here \(\epsilon>0\) gives
\(
|\epsilon s(u)-\epsilon s(v)|
=\epsilon|s(u)-s(v)|\leq\epsilon\delta
\).
% lean: CausalSmith.Stat.OptvalueVanishingoverlapRate.armwise_modulus_extension

\emph{2. The zero-vector case.}
Suppose \(u=0\). Nonnegativity and the definition of the denominator imply
\[
0\leq v_{a1}\leq s_a(v)\leq D_a(v),
\qquad
0\leq\frac{s(v)v_{a1}}{D_a(v)}\leq s(v).
\]
Taking the maximum over the two arms gives
\[
0\leq F_\epsilon(v)\leq s(v).
\]
Also,
\[
s(v)=|s(u)-s(v)|\leq\delta,
\qquad
\sum_{a=0}^1\frac{s(v)}{D_a(v)}\Delta_a\geq0.
\]
Since \(F_\epsilon(u)=0\), it follows that
\[
|F_\epsilon(u)-F_\epsilon(v)|
=F_\epsilon(v)
\leq\delta
\leq2\delta+
2\sum_{a=0}^1\frac{s(v)}{D_a(v)}\Delta_a.
\]
% lean: CausalSmith.Stat.OptvalueVanishingoverlapRate.armwiseExtension_nonneg_le_totalMassVec

\emph{3. A coordinate-ratio estimate.}
Now suppose \(u\neq0\). The same nonnegative-coordinate argument gives \(s(u)>0\), and hence, for each arm,
\[
D_a(u)\geq\epsilon s(u)>0,\qquad
0\leq u_{a1}\leq D_a(u),\qquad
0\leq v_{a1}\leq D_a(v).
\]
We first establish
\[
\left|\frac{u_{a1}}{D_a(u)}-\frac{v_{a1}}{D_a(v)}\right|
\leq
\frac{|u_{a1}-v_{a1}|+|D_a(u)-D_a(v)|}{D_a(v)}.
\]
If \(D_a(u)\leq D_a(v)\), the identity
\[
D_a(v)u_{a1}-D_a(u)v_{a1}
=
D_a(u)(u_{a1}-v_{a1})
+(D_a(v)-D_a(u))u_{a1}
\]
and \(0\leq u_{a1}\leq D_a(u)\) give
\[
\begin{aligned}
|D_a(v)u_{a1}-D_a(u)v_{a1}|
&\leq D_a(u)|u_{a1}-v_{a1}|
 +(D_a(v)-D_a(u))u_{a1}\\
&\leq D_a(u)
 \bigl(|u_{a1}-v_{a1}|+|D_a(u)-D_a(v)|\bigr).
\end{aligned}
\]
If \(D_a(v)\leq D_a(u)\), use instead
\[
D_a(v)u_{a1}-D_a(u)v_{a1}
=
D_a(v)(u_{a1}-v_{a1})
-(D_a(u)-D_a(v))v_{a1}.
\]
Since \(0\leq v_{a1}\leq D_a(v)\leq D_a(u)\),
\[
\begin{aligned}
|D_a(v)u_{a1}-D_a(u)v_{a1}|
&\leq D_a(v)|u_{a1}-v_{a1}|
 +(D_a(u)-D_a(v))v_{a1}\\
&\leq D_a(u)
 \bigl(|u_{a1}-v_{a1}|+|D_a(u)-D_a(v)|\bigr).
\end{aligned}
\]
In both cases, divide by the positive product \(D_a(u)D_a(v)\), using
\[
\frac{u_{a1}}{D_a(u)}-\frac{v_{a1}}{D_a(v)}
=
\frac{D_a(v)u_{a1}-D_a(u)v_{a1}}{D_a(u)D_a(v)},
\]
to obtain the coordinate-ratio estimate.
% lean: CausalSmith.Stat.OptvalueVanishingoverlapRate.armwise_ratio_bound

\emph{4. Scaling the ratios and comparing the maxima.}
For the two nonzero vectors under consideration, define
\[
g_a(\xi):=\frac{s(\xi)\xi_{a1}}{D_a(\xi)},
\qquad \xi\in\{u,v\},\quad a\in\{0,1\}.
\]
The identity
\[
g_a(u)-g_a(v)
=
(s(u)-s(v))\frac{u_{a1}}{D_a(u)}
+s(v)\left(\frac{u_{a1}}{D_a(u)}
           -\frac{v_{a1}}{D_a(v)}\right)
\]
gives
\[
\begin{aligned}
|g_a(u)-g_a(v)|
&\leq |s(u)-s(v)|\frac{u_{a1}}{D_a(u)}
+s(v)\left|\frac{u_{a1}}{D_a(u)}
           -\frac{v_{a1}}{D_a(v)}\right|\\
&\leq\delta+
\frac{s(v)}{D_a(v)}
\bigl(\Delta_a+|D_a(u)-D_a(v)|\bigr).
\end{aligned}
\]
Indeed, \(0\leq u_{a1}/D_a(u)\leq1\), and
\(
|u_{a1}-v_{a1}|\leq\Delta_a
\).
% lean: CausalSmith.Stat.OptvalueVanishingoverlapRate.armwise_coordinate_ratio_bound
Moreover,
\[
\epsilon\frac{s(v)}{D_a(v)}
=\frac{\epsilon s(v)}{D_a(v)}\leq1.
\]
Substituting the denominator perturbation bound therefore yields
\[
\begin{aligned}
|g_a(u)-g_a(v)|
&\leq\delta+
\frac{s(v)}{D_a(v)}(2\Delta_a+\epsilon\delta)\\
&\leq2\delta+2\frac{s(v)}{D_a(v)}\Delta_a.
\end{aligned}
\]
Finally, comparison of the two entries of the maximum gives
\[
\begin{aligned}
|F_\epsilon(u)-F_\epsilon(v)|
&=|\max\{g_0(u),g_1(u)\}-\max\{g_0(v),g_1(v)\}|\\
&\leq\max_{a\in\{0,1\}}|g_a(u)-g_a(v)|\\
&\leq\max_{a\in\{0,1\}}
 \left\{2\delta+2\frac{s(v)}{D_a(v)}\Delta_a\right\}\\
&\leq2\delta+
2\sum_{a=0}^1\frac{s(v)}{D_a(v)}\Delta_a.
\end{aligned}
\]
The last inequality uses the nonnegativity of both weighted perturbations. This proves the first assertion.
% lean: CausalSmith.Stat.OptvalueVanishingoverlapRate.armwise_modulus_extension

\emph{5. Observed arm-mass floors.}
For the second assertion, fix the stated law and cell, and take \(v=q_x\). The atom probabilities are nonnegative. The definitions in \cref{obj:def:observed-class,obj:def:armwise-extension}, together with \(p_x>0\), give
\[
p_x=s_0(v)+s_1(v),\qquad
p_x\pi_x=s_1(v),\qquad
p_x(1-\pi_x)=s_0(v).
\]
The last identity follows by subtracting the second from the first. Membership in the observed class supplies \cref{obj:ass:shrinking-overlap}, so
\[
\epsilon\leq\pi_x,\qquad \epsilon\leq1-\pi_x.
\]
Multiplication by \(p_x>0\) gives, separately for the two arms,
\[
\epsilon p_x\leq p_x(1-\pi_x)=s_0(v),
\qquad
\epsilon p_x\leq p_x\pi_x=s_1(v).
\]
Consequently,
\[
s_a(v)\geq\epsilon p_x>0,\qquad a\in\{0,1\}.
\]
% lean: CausalSmith.Stat.OptvalueVanishingoverlapRate.observed_armwise_sqrt_bound

\emph{6. The weighted square-root estimate.}
For each arm, nonnegativity of its two coordinates gives
\[
\begin{aligned}
2s_a(v)-(\sqrt{v_{a0}}+\sqrt{v_{a1}})^2
&=(\sqrt{v_{a0}}-\sqrt{v_{a1}})^2\geq0,\\
(\sqrt{v_{a0}}+\sqrt{v_{a1}})^2
&\leq2s_a(v).
\end{aligned}
\]
Both sides of the resulting square-root comparison are nonnegative, hence
\[
\sqrt{v_{a0}}+\sqrt{v_{a1}}\leq\sqrt{2s_a(v)}.
\]
% lean: CausalSmith.Stat.OptvalueVanishingoverlapRate.sqrt_pair_le_sqrt_twice_sum
Multiplying the squared bound by \(p_x^2\), and then using the arm-mass floor, yields
\[
\begin{aligned}
\bigl(p_x(\sqrt{v_{a0}}+\sqrt{v_{a1}})\bigr)^2
&\leq2p_x^2s_a(v)\\
&=\frac{2p_xs_a(v)}{\epsilon}(\epsilon p_x)\\
&\leq\frac{2p_xs_a(v)}{\epsilon}s_a(v)\\
&=\frac{2p_x}{\epsilon}s_a(v)^2\\
&=\left(\sqrt{\frac{2p_x}{\epsilon}}\,s_a(v)\right)^2.
\end{aligned}
\]
The compared quantities are nonnegative. Taking square roots and dividing by \(s_a(v)>0\) proves
\[
\frac{p_x}{s_a(v)}
\bigl(\sqrt{v_{a0}}+\sqrt{v_{a1}}\bigr)
\leq\sqrt{\frac{2p_x}{\epsilon}},
\qquad a\in\{0,1\}.
\]
% lean: CausalSmith.Stat.OptvalueVanishingoverlapRate.weighted_sqrt_pair_bound

\emph{7. Summing the two arms.}
Adding the two armwise bounds and factoring the square root gives
\[
\begin{aligned}
\sum_{a=0}^1\frac{p_x}{s_a(v)}
\sum_{y=0}^1\sqrt{v_{ay}}
&\leq2\sqrt{\frac{2p_x}{\epsilon}}\\
&=2\sqrt{2}\sqrt{\frac{p_x}{\epsilon}}.
\end{aligned}
\]
This proves the second assertion and completes the proof.
% lean: CausalSmith.Stat.OptvalueVanishingoverlapRate.armwise_modulus
\end{proof}
